\documentclass[11pt]{article}
\usepackage[margin=1in]{geometry}
\usepackage{amsmath,amssymb,amsthm}
\usepackage{graphicx}
\usepackage{hyperref}

\newtheorem{theorem}{Theorem}

\newcommand{\C}{\mathbb C}
\newcommand{\N}{\mathbb N}
\newcommand{\ellTwo}{\ell_2}
\newcommand{\cnull}{c_{00}}
\newcommand{\LL}{\mathcal L}
\newcommand{\J}{J}
\newcommand{\norm}[1]{\left\lVert #1 \right\rVert}

\title{A nilpotent right-perturbation criterion for hypercyclic generalized derivations}
\author{Run fa\_banach\_001, agent\_lane\_10}
\date{2026-06-21}

\begin{document}
\maketitle

\section*{Status}

This is a likely valid full positive answer to Section 6, Question 1 of
arXiv:1605.07409,
Clifford Gilmore, \emph{Dynamics of Generalised Derivations and Elementary
Operators}.  It gives a broad positive answer to the first question in Section
6 of that paper, which asks whether reasonable sufficient conditions exist on
the pair \((A,B)\).  It is not a characterization of all hypercyclic
generalized derivations.

\section*{Source question}

The source paper defines, for operators \(A,B\in \LL(X)\),
\[
  \tau_{A,B}(S)=L_A(S)-R_B(S)=AS-SB.
\]
In Section 6, Question 1 asks:

\begin{quote}
Do reasonable sufficient conditions exist on the pair \((A,B)\) that induce
hypercyclic \(\tau_{A,B}\) on separable Banach ideals?
\end{quote}

The following crop shows the question in the source PDF.

\begin{center}
\includegraphics[width=\textwidth]{figures/open_problem_crop.png}
\end{center}

\section*{Idea of the proof}

The source paper already shows how to turn one-sided hypercyclicity into
hypercyclic generalized derivations when the second operator is a scalar
multiple of the identity.  The point here is that scalarity can be replaced by
nilpotence.  Suppose \(A\) has a right inverse \(F\) with spectral radius less
than one, and suppose a dense subspace is eventually killed by powers of \(A\).
If \(B=N\) is nilpotent, then \(R_N\) is a commuting nilpotent perturbation of
the left multiplier \(L_A\).  The expression \((L_A-R_N)^nL_F^n\) is an
invertible polynomial in the nilpotent operator \(L_FR_N\).  Its inverse grows
at most polynomially in \(n\), while \(L_F^n\) decays exponentially.  This
gives explicit approximate inverse maps for the Hypercyclicity Criterion.  On
finite-rank operators whose ranges lie in the eventually killed dense subspace,
the powers of \(L_A-R_N\) are eventually zero.

\section*{Theorem}

Let \(X\) be a Banach space, let \(\J\subset \LL(X)\) be a separable admissible
Banach ideal in the sense used in the source paper, and let \(A,N\in\LL(X)\).
Assume that:
\begin{enumerate}
\item \(N\neq 0\) is nilpotent;
\item there is \(F\in\LL(X)\) such that \(AF=I_X\) and \(r(F)<1\);
\item there is a dense linear subspace \(X_0\subset X\) such that, for every
  \(x\in X_0\), \(A^n x=0\) for all sufficiently large \(n\).
\end{enumerate}
Then
\[
  \tau_{A,N}=L_A-R_N:\J\to\J
\]
satisfies the Hypercyclicity Criterion.  In particular, it is hypercyclic.

Thus every pair \((A,N)\) satisfying the three conditions above is a concrete
sufficient condition for the first final question of the source paper, with
\(N\) nonzero and non-scalar.  Since that question asks whether such reasonable
sufficient conditions exist, this is a positive answer to the question as
stated.

\section*{Proof}

Let \(m\geq 2\) be such that \(N^m=0\).  Then \(R_N^m=0\), and all left
multipliers commute with all right multipliers.  Since \(r(F)<1\), there are
constants \(C>0\) and \(0<\rho<1\) such that
\[
  \norm{F^n}\leq C\rho^n\qquad(n\geq 0).
\]

Set
\[
  T=L_A-R_N \quad\hbox{on }\J,
  \qquad U=L_F R_N .
\]
For every \(n\geq m\), the binomial formula and \(R_N^m=0\) give
\[
  T^n L_F^n
   = \sum_{j=0}^{m-1}(-1)^j \binom{n}{j}
        L_{A^{\,n-j}F^n} R_N^j
   = \sum_{j=0}^{m-1}(-1)^j \binom{n}{j} U^j
   =: Q_n(U).
\]
Here we used \(A^{\,n-j}F^n=F^j\), which follows from \(AF=I_X\).  Since
\(U^m=0\) and \(Q_n(0)=1\),
the operator \(Q_n(U)\) is invertible, with inverse
\[
  Q_n(U)^{-1}
  = \sum_{r=0}^{m-1}\left(I-Q_n(U)\right)^r .
\]
The norm of \(Q_n(U)^{-1}\) grows at most polynomially in \(n\), because it is
a polynomial of degree \(m-1\) in the nilpotent \(U\), with coefficients
polynomially bounded in the binomial coefficients \(\binom{n}{j}\).

Define
\[
  S_n=L_F^n Q_n(U)^{-1}\qquad(n\geq m).
\]
Then
\[
  T^n S_n = T^n L_F^n Q_n(U)^{-1}=I_\J.
\]
Moreover,
\[
  \norm{S_n}\leq \norm{L_F^n}\,\norm{Q_n(U)^{-1}}
       \leq C\rho^n\,\norm{Q_n(U)^{-1}}\longrightarrow 0,
\]
since exponential decay dominates the polynomial growth.  Hence
\[
  S_n Y\to 0
  \quad\hbox{and}\quad
  T^n S_n Y=Y
\]
for every \(Y\in \J\).

It remains to find a dense subspace on which \(T^n\to 0\).  Let
\[
  \mathcal D=\operatorname{span}\{x^*\otimes x:\ x^*\in X^*,\ x\in X_0\}.
\]
The admissibility assumptions used in the source paper imply that finite-rank
operators are dense in \(\J\); since \(X_0\) is dense in \(X\), \(\mathcal D\)
is dense in \(\J\), by the same rank-one approximation argument used in the
source paper.  If \(Y\in \mathcal D\), then the range of \(Y\) is contained in
a finite-dimensional subspace spanned by finitely many vectors of \(X_0\).
For each fixed \(0\leq j<m\), the operator \(R_N^jY=YN^j\) has range contained
in the same finite-dimensional subspace.  By the eventual-killing hypothesis on
\(X_0\), there is \(M_Y\) such that
\[
  L_{A^{\,n-j}}R_N^jY=0
\]
whenever \(n-j\geq M_Y\).  Using the binomial expansion
\[
  T^nY=\sum_{j=0}^{m-1}(-1)^j \binom{n}{j}
        L_{A^{\,n-j}}R_N^jY,
\]
we get \(T^nY=0\) for all sufficiently large \(n\).  Thus \(T^nY\to 0\) for
every \(Y\in \mathcal D\).

The dense subspace \(\mathcal D\), the sequence \(n\mapsto n\), and the maps
\(S_n\) verify the Hypercyclicity Criterion for \(T=\tau_{A,N}\).  Therefore
\(\tau_{A,N}\) is hypercyclic on \(\J\).

\section*{Concrete example}

Let \(X=\ell_2\).  Let \(W_\lambda\) be \(\lambda\) times the backward shift
for some \(\lambda>1\), and let \(F=\lambda^{-1}\) times the forward shift.
Then \(W_\lambda F=I_X\), \(r(F)=\lambda^{-1}<1\), and \(W_\lambda^n x=0\)
eventually for every \(x\in c_{00}\).  Taking
\[
  N=e_0\otimes e_1^*,
\]
so that \(Ne_1=e_0\), \(Ne_k=0\) for \(k\neq 1\), and \(N^2=0\), the theorem
gives a hypercyclic generalized derivation \(\tau_{W_\lambda,N}\) on every
separable admissible Banach ideal \(\J\subset\LL(\ell_2)\).

\section*{Verification, novelty check, and limitations}

The proof uses only the source paper's admissible-ideal framework and the
standard Hypercyclicity Criterion.  No computational verification is needed.

This packet gives a broad sufficient-condition theorem, not a classification of
all hypercyclic pairs \((A,B)\).  It also does not address the second final
question in the source paper about contrasting hypercyclic behavior across
different ideals.

Before promotion to a full packet, the run's cheap indexes and local parsed
source corpus were searched for ``1605.07409'', ``generalised derivations'',
``generalized derivations'', ``elementary operators'', ``nilpotent right
perturbation'', ``\(L_A-R_N\)'', and related Hypercyclicity Criterion phrases.
Targeted web/arXiv searches on 2026-06-21 for the exact source-question phrase
and close variants such as ``nilpotent right perturbation generalized
derivation hypercyclic'' found the source paper but no later matching
criterion.  The novelty confidence is therefore bounded, but no duplicate was
found.

\section*{References}

\begin{thebibliography}{9}

\bibitem{Gilmore}
C. Gilmore,
\emph{Dynamics of Generalised Derivations and Elementary Operators},
arXiv:1605.07409.

\end{thebibliography}

\end{document}
